% =====================
\section{資源・調達管理}

担当B/Cが使用する部品・機材を表\ref{tab:resources}に示す．

\begin{table}[H]
  \centering
  \begin{tabularx}{\linewidth}{lXll}
    \toprule
    部品・機材 & 用途 & 数量 & 調達方法 \\
    \midrule
    Arduino Uno R4 WiFi &
      子機（楽器ユニット）制御 & 5台 & 学科貸出 \\
    赤外線距離センサー &
      腕位置の連続検出（型番は試作時に確定） & 5個 & 学科貸出または購入 \\
    フォトトランジスタ &
      腕先端LED検出 & 5個 & 学科貸出または購入 \\
    抵抗・コンデンサ等 &
      センサー回路用受動部品 & 適量 & 学科貸出 \\
    USBケーブル（Type-B） &
      Arduino--PCシリアル接続 & 5本 & 学科貸出 \\
    PC（Mac） &
      Processing実行・シリアル受信 & 5台 & チームメンバー持参 \\
    \bottomrule
  \end{tabularx}
  \caption{資源・調達一覧（担当B/C関連）}
  \label{tab:resources}
\end{table}
